% ---------------------------------------------------------
% cv_template.tex.jinja — Master layout for the full CV
% ---------------------------------------------------------

% ---------------------------------------------------------
% _preamble.tex.j2 — Sets up LaTeX document, packages, etc.
% ---------------------------------------------------------

% Use standard article class for simple documents
\documentclass[letterpaper, 11pt]{article}

% Language and symbols
\usepackage[english]{babel}
\usepackage{latexsym}

% Full-page usage with no headers/footers
\usepackage[empty]{fullpage}

% Section formatting and symbols
\usepackage{titlesec}
\usepackage{marvosym}

% Colors and font styling
\usepackage[usenames, dvipsnames]{color}
\usepackage{enumitem}

% URL handling
\usepackage[hidelinks]{hyperref}

% Tables, icons, and typography
\usepackage{tabularx}
\usepackage{fontawesome5}
\usepackage{multicol}
\usepackage{graphicx}
\usepackage{CormorantGaramond}
\usepackage{charter}
\usepackage{tikz}
\usepackage{graphicx}

% Unicode support for PDF output
\input{glyphtounicode}
\pdfgentounicode=1

% Page layout tweaks
\addtolength{\oddsidemargin}{-0.6in}
\addtolength{\evensidemargin}{-0.5in}
\addtolength{\textwidth}{1.19in}
\addtolength{\topmargin}{-.7in}
\addtolength{\textheight}{1.4in}
\raggedbottom
\raggedright
\setlength{\tabcolsep}{0in}

% Hyperlink color styling
\hypersetup{
    colorlinks=true,
    linkcolor=black,
    urlcolor=linki,
}
\urlstyle{same}

\setlength{\multicolsep}{-3.0pt}
\setlength{\columnsep}{-1pt}
% ---------------------------------------------------------
% _colors.tex.j2 — Custom color definitions and assignments
% ---------------------------------------------------------

% Define custom colors using HTML hex or RGB
\definecolor{cvblue}{HTML}{0E5484}
\definecolor{black}{HTML}{130810}
\definecolor{darkcolor}{HTML}{0F4539}
\definecolor{andreblue}{RGB}{31, 119, 180}
\definecolor{linki}{HTML}{2020DF}

% Semantic color assignments — tweakable if theme changes
\colorlet{name}{black}
\colorlet{tagline}{darkcolor}
\colorlet{heading}{darkcolor}
\colorlet{headingrule}{cvblue}
\colorlet{accent}{darkcolor}

%========================================================================================
% _macros.tex
%
% This file defines all the custom LaTeX commands (macros) used to structure and
% style the resume.
%========================================================================================


%----------------------------------------------------------------------------------------
%	SECTION FORMATTING
%----------------------------------------------------------------------------------------
% Defines the style for \section headings using the `titlesec` package.
\titleformat{\section}{
  \vspace{-4pt}\scshape\raggedright\large\bfseries
}{}{0em}{}[\color{black}\titlerule \vspace{-5pt}]


%----------------------------------------------------------------------------------------
%	HEADING & ENTRY COMMANDS
%----------------------------------------------------------------------------------------

% Defines a wrapper environment for a list of entries (e.g., jobs, projects).
% It creates a clean vertical list with no bullets or indentation.
\newcommand{\resumeSubHeadingListStart}{\begin{itemize}[leftmargin=0.0in, label={}]}
\newcommand{\resumeSubHeadingListEnd}{\end{itemize}}

% Creates a standard entry for a job, project, or education section.
% ARGUMENTS: {Title}{Location}{Subtitle/Company}{Dates}{Description list}
\newcommand{\resumeSubheading}[5]{
  \vspace{-2pt}\item
    \begin{tabular*}{1.0\textwidth}[t]{@{}p{0.85\textwidth}@{\extracolsep{\fill}}r@{}}
      \textbf{\normalsize #1} & \textbf{\small #2} \\
      \textit{\small #3} & \textit{\small #4} \\
      \small\hspace{1em}\begin{minipage}[t]{\linewidth}#5\end{minipage} & \\
    \end{tabular*}\vspace{-7pt}
}

% Creates a simpler subheading for sections like Hobbies or Activities.
% ARGUMENTS: {Title}{Description list}
\newcommand{\hobbySubheading}[2]{
  \vspace{3pt}\item
    \begin{tabular*}{1.0\textwidth}[t]{@{}p{\textwidth}@{\extracolsep{\fill}}r@{}}
      \textbf{\normalsize #1} &  \\
      \small\hspace{1em}\begin{minipage}[t]{\linewidth}#2\end{minipage} &\\
    \end{tabular*}\vspace{-5pt}
}


%----------------------------------------------------------------------------------------
%	BULLET POINT LIST COMMANDS
%----------------------------------------------------------------------------------------

% Defines a standard, compact, indented list for bullet points.
% Spacing and bullet style are defined here for consistency.
\newcommand{\resumeItemListStart}{
  \begin{itemize}[leftmargin=0.15in, label=\tiny$\bullet$, topsep=0pt, itemsep=0em, partopsep=0pt, parsep=0pt]
}
\newcommand{\resumeItemListEnd}{\end{itemize}}

% Creates an alias for the hobby list. It will now use the exact same styling
% as the standard resume item list, ensuring visual consistency.
\let\hobbyItemListStart\resumeItemListStart
\let\hobbyItemListEnd\resumeItemListEnd


%----------------------------------------------------------------------------------------
%	ITEM-LEVEL COMMANDS
%----------------------------------------------------------------------------------------

% A standard single-column list item.
% ARGUMENTS: {Description}
\newcommand{\resumeItem}[1]{
  \item\small{#1}
}

% A two-column list item, with a layout consistent with \resumeSubheading.
% It correctly handles multi-line descriptions.
% ARGUMENTS: {Description}{Right-aligned text (e.g., Date)}
\newcommand{\hobbyItem}[2]{
  \item
    \begin{tabular*}{\linewidth}[t]{@{}p{0.85\linewidth}@{\extracolsep{\fill}}r@{}}
      #1 & \textit{#2}
    \end{tabular*}
}

\begin{document}


% ---------------------------------------------------------
% _header.tex.j2 - Renders contact info line below the name
% ---------------------------------------------------------

\begin{center}

{\huge\textbf{\textsc{ Marie Sklodowska Curie }}} \\

        \vspace{5pt}
% --- Phone (Plain Text) ---
+33 1 40 46 75 00
% --- Email (Special mailto link) ---
~--~ \href{mailto:marie.curie@sorbonne.fr}{marie.curie@sorbonne.fr}
% --- Location (Plain Text) ---
~--~ Paris, France
% --- GitHub (Link) ---
~--~ \href{https://www.nobelprize.org/prizes/physics/1903/marie-curie/biographical/}{Nobel Prize Profile }
% --- LinkedIn (Link) ---
\vspace{-7pt}
\end{center}
% ---------------------------------------------------------
% _education.tex.j2 — Renders the Education section
% ---------------------------------------------------------

\section{\texorpdfstring{\color{andreblue}Education}{Education}}
\resumeSubHeadingListStart

  \resumeSubheading
    {Doctor of Science in Physics \textnormal{\small~Awarded 'tres honorable' (very honorable)}}
    {Paris, FR}
    {University of Paris (Sorbonne)}
    {Jun 1903}
{\resumeItemListStart
        \resumeItem{Thesis established the new scientific field of radioactivity, hypothesizing that radiation was an atomic property.}
        \resumeItem{Presented novel methodologies for the measurement of radioactivity, forming a cornerstone for future research in atomic physics.}
    \resumeItemListEnd
    }

  \resumeSubheading
    {Degree in Mathematics \textnormal{\small~Graduated second in class}}
    {Paris, FR}
    {University of Paris (Sorbonne)}
    {Jul 1894}
{\resumeItemListStart
        \resumeItem{Acquired advanced mathematical knowledge essential for theoretical physics and quantitative modeling.}
        \resumeItem{Pursued as a second degree to build a multidisciplinary foundation for scientific research.}
    \resumeItemListEnd
    }

  \resumeSubheading
    {Degree in Physics \textnormal{\small~Graduated first in class}}
    {Paris, FR}
    {University of Paris (Sorbonne)}
    {Jul 1893}
{\resumeItemListStart
        \resumeItem{Achieved top academic standing in a rigorous curriculum covering mechanics, electricity, and thermodynamics.}
        \resumeItem{Built a comprehensive understanding of classical physics, enabling the identification of phenomena that defied existing theories.}
    \resumeItemListEnd
    }

  \resumeSubheading
    {Informal Scientific Training \textnormal{\small~Completed Clandestine Studies}}
    {Warsaw, PL}
    {Flying University (Uniwersytet Latajacy)}
    {1886 - 1889}
{\resumeItemListStart
        \resumeItem{Participated in a clandestine underground university to pursue scientific education forbidden to women in Poland.}
        \resumeItem{Gained foundational knowledge in chemistry and physics, sparking an interest in empirical research.}
    \resumeItemListEnd
    }


\resumeSubHeadingListEnd
\vspace{-16pt}
% ---------------------------------------------------------
% _experience.tex.j2 — Professional experience section
% ---------------------------------------------------------

\section{\texorpdfstring{\color{andreblue}Experience}{Experience}}
\resumeSubHeadingListStart

  \resumeSubheading
    { \textbf{Director of the Radium Institute} }
    { \textbf{\small Paris, FR} }
    { \textit{\normalsize at University of Paris} }
    { \textit{\small 1914 - 1934} }
    {
      \resumeItemListStart
          \resumeItem{ Held the post of Professor of General Physics, succeeding Pierre Curie and becoming the first female professor at the Sorbonne (appointed 1906). }
          \resumeItem{ Directed a world-leading research center for chemistry, physics, and medicine, mentoring numerous doctoral students. }
          \resumeItem{ Established a new scientific discipline and oversaw fundamental research into the nature of radioactivity and its applications. }
      \resumeItemListEnd
    }
  \resumeSubheading
    { \textbf{Director of the Red Cross Radiology Service} }
    { \textbf{\small France} }
    { \textit{\normalsize at French Red Cross} }
    { \textit{\small 1914 - 1918} }
    {
      \resumeItemListStart
          \resumeItem{ Conceptualized and developed mobile radiography units ('petites Curies') to diagnose injuries at the front lines. }
          \resumeItem{ Authored training materials and personally instructed 150 women as the first military radiology technicians. }
      \resumeItemListEnd
    }
  \resumeSubheading
    { \textbf{Researcher \& Scientist} }
    { \textbf{\small Paris, FR} }
    { \textit{\normalsize at School of Chemistry and Physics (ESPCI)} }
    { \textit{\small 1897 - 1906} }
    {
      \resumeItemListStart
          \resumeItem{ Conducted foundational research on 'Becquerel rays', coining the term 'radioactivity'. }
          \resumeItem{ Discovered and scientifically described two new elements in the periodic table: polonium (Po, 1898) and radium (Ra, 1898). }
          \resumeItem{ Determined the atomic weight of radium, solidifying its status as a new element, and investigated its physical and chemical properties. }
      \resumeItemListEnd
    }

\resumeSubHeadingListEnd

% ---------------------------------------------------------
% _projects.tex.j2 — Lists personal or academic projects
% ---------------------------------------------------------



% ---------------------------------------------------------
% _skills.tex.j2 — Renders technical skills, tools, or language proficiencies
% ---------------------------------------------------------

\section{\texorpdfstring{\color{andreblue}Skills}{Skills}}

\resumeSubHeadingListStart

  \resumeItem{\textbf{Scientific Expertise:} Radioactivity, Nuclear Physics (emerging), Analytical Chemistry, Crystallography, Electrometry}
  \vspace{-7pt}
  \resumeItem{\textbf{Laboratory Techniques:} Fractional Crystallization, Chemical Separation \& Isolation, Electrometric Measurement, Radiography (X-ray), Glassblowing, Sample Preparation}
  \vspace{-7pt}
  \resumeItem{\textbf{Technical Instrumentation:} Piezoelectric Electrometer, Mobile X-ray Units ('petites Curies'), Radiation Shielding Protocols, Calorimeters for heat measurement}
  \vspace{-7pt}
  \resumeItem{\textbf{Languages:} Native in Polish, Fluent in French, Proficient in German \& Russian, Reading knowledge of English}
  \vspace{-7pt}

\resumeSubHeadingListEnd

% ---------------------------------------------------------
% _extracurricular.tex.j2 — Renders extracurricular activities, leadership, and interests
% ---------------------------------------------------------

\section{\texorpdfstring{\color{andreblue}Extracurricular}{Extracurricular}}

\resumeSubHeadingListStart


  \hobbySubheading
    {Scientific Leadership \& Outreach}
    {
      % Start the list of items for this group.
      \hobbyItemListStart
          \hobbyItem{Active participant and only female member of the prestigious Solvay Conferences on Physics and Chemistry}{1911 - 1933}
          \hobbyItem{Founded and directed the Radium Institute in her native Warsaw (now the Maria Sklodowska-Curie National Research Institute of Oncology)}{1925 - 1932}
          \hobbyItem{Maintained extensive correspondence on theoretical and experimental physics with leading scientists, including Albert Einstein, Ernest Rutherford, and Max Planck}{c. 1900 - 1934}
      \hobbyItemListEnd
    }

  \hobbySubheading
    {Hobbies \& Personal Interests}
    {
      % Start the list of items for this group.
      \hobbyItemListStart
          \hobbyItem{Enjoyed long bicycle trips and extensive travel, including tours of the United States to fund her research}{Lifelong}
          \hobbyItem{Kept meticulous and detailed personal and scientific journals throughout her life}{c. 1885 - 1934}
          \hobbyItem{Swimming and hiking in the mountains}{Lifelong}
      \hobbyItemListEnd
    }

\resumeSubHeadingListEnd

    % Push this content to the bottom of the page,
    % regardless of how much text precedes it.
    \vspace*{\fill}%

    \makebox[2.5in]{\hrulefill}\\
    Paris\hspace{80pt} \today
    \\
    \vspace{-51pt}
    \hspace{20pt}\includegraphics[width=5cm]{assets/SignatureMarieCurie.pdf}


\end{document}